\section{Restrictions on Picard kernels}
\label{sec:picard-kernel-constraints}
We examine two ways in which a torsion divisor class might disappear on a
smooth ample curve. A purely inseparable multisection of a quasi-elliptic
fibration has the desired effect on torsion, but the intersection calculation
below prevents its ampleness over a base of positive genus. For an arbitrary
order-$p$ class, normalization of its torsor gives a different constraint:
the lifted divisor is reduced, and the conductor along it is determined
exactly. In the equality case, a group action extends to the normalized
cover, and a zero scheme on the smooth base counts its isolated
singularities. These constraints leave the general point-kernel
question open.

For the direct route to the genus results, first read the normalized
torsor calculation in Proposition~\ref{prop:normalized-mup-conductor}.
The adjoint argument is in Section~\ref{subsec:canonical-subsystem-strictness},
and the very ample equality structure is in
Section~\ref{subsec:conductor-equality-structure}; these prove
Theorems~\ref{thm:strict-picard-all} and
\ref{thm:conductor-equality-structure}.
Section~\ref{subsec:positive-square-picard} gives a self-contained
route to the quantitative positive-square bound of
Theorem~\ref{thm:positive-excess-picard}.
The intervening ruled, K3 and rational-cover arguments treat the
merely ample cases, whose equality configurations have separate
hypotheses and remain relevant without very ampleness.

\begin{proposition}\label{prop:quasielliptic-multisection}
Let $k$ be algebraically closed of characteristic $p=2$ or $3$, and let
$f:S\to C$ be a quasi-elliptic fibration from a smooth projective integral
surface to a smooth projective integral curve. Thus $f_*\OO_S=\OO_C$ and
the generic fibre is a regular geometrically integral curve of arithmetic
genus one that is not smooth. Let $R\subset S$ be a smooth integral
horizontal curve for which $R\to C$ is purely inseparable of degree
$q=p^e\geq2$. If $g=g(C)$, then
\[
 R^2\leq\left(2-\frac{5q}{3}\right)(g-1).
\]
In particular, $R$ is not ample when $g\geq1$. The fibration need not be
relatively minimal.
\end{proposition}
\begin{proof}
First $g(R)=g$. Indeed, put $L=k(R)$ and $K=k(C)$. A one-variable
function field over a perfect field satisfies $[L:L^{p^e}]=p^e$.
The purely inseparable extension $L/K$ has exponent at most $e$, so
$L^{p^e}\subset K$; equality of degrees gives $K=L^{p^e}$.
The map of smooth curves is therefore an iterated relative Frobenius
followed by an isomorphism, and Frobenius twisting preserves genus.

Contract the vertical exceptional curves to obtain a relatively minimal
smooth model $f_0:S_0\to C$, with contraction $\beta:S\to S_0$.
Write $\mathcal E$ for the line-bundle quotient of
$R^1(f_0)_*\OO_{S_0}$ by its torsion subsheaf $\mathcal T$, and put
$\mathcal L=\mathcal E^\vee$ and $t=\operatorname{length}(\mathcal T)$.
The canonical-bundle formula, including wild multiple fibres, is
\[
 K_{S_0}\sim f_0^*(K_C+\mathcal L)+\sum_i a_iP_i,
 \qquad 0\leq a_i<m_i,\qquad f_0^*(c_i)=m_iP_i;
\]
see \cite[Theorem~5.6, author pp.~23--24]{Liedtke} and
\cite[Theorem~4.1.6, pp.~364--365]{CDL}.
Leray and Riemann--Roch on $C$ give
\[
 \chi(\OO_{S_0})
 =\chi(\OO_C)-\chi(\mathcal E)-t
 =-\deg\mathcal E-t,
 \qquad \deg\mathcal L=\chi(\OO_{S_0})+t.
\]
Every term in the canonical formula is numerically a rational multiple
of a fibre, so $K_{S_0}^2=0$.

The quasi-elliptic Euler formula is
\[
 c_2(S_0)=4(1-g)+\sum_{c\in C}\bigl(b_2((S_0)_c)-1\bigr)
 \geq4(1-g);
\]
here $b_2$ is the second $\ell$-adic Betti number for $\ell\ne p$,
and only finitely many summands are nonzero
\cite[Corollary~4.1.12 and the preceding quasi-elliptic discussion,
p.~369]{CDL}. Each fibre has at least one irreducible component,
so each summand is nonnegative. Noether's formula now yields
$\chi(\OO_{S_0})=c_2(S_0)/12\geq(1-g)/3$.

Let $R_0=\beta(R)$ with its reduced structure. It is horizontal of degree
$q$ over $C$. Its arithmetic genus need not equal $g$, so we apply
adjunction on $S$, where $R$ is smooth. The divisor
$K_S-\beta^*K_{S_0}$ is effective and exceptional, with no $R$ component.
Its intersection with $R$ is nonnegative. The effective vertical terms
in the canonical formula also have nonnegative intersection with $R_0$.
Consequently
\[
 K_S\cdot R\geq K_{S_0}\cdot R_0
 \geq q\bigl(2g-2+\chi(\OO_{S_0})+t\bigr).
\]
Adjunction gives
\[
 \begin{split}
 R^2
 &=2g-2-K_S\cdot R\\
 &\leq(1-q)(2g-2)-q\chi(\OO_{S_0})-qt\\
 &\leq\left(2-\frac{5q}{3}\right)(g-1).
 \end{split}
\]
This is nonpositive for $g\geq1$ and $q\geq2$, whereas an ample curve
on a projective surface has positive self-intersection.
\end{proof}

A nontrivial $p$-torsion line bundle on $C$ requires $g(C)\geq1$.
Pullback by an iterated relative Frobenius kills such a class, after the
corresponding Frobenius twist of the curve. The proposition explains why
this mechanism on a quasi-elliptic fibration cannot supply a smooth ample
multisection, even after point blowups. It does not exclude maps with an
additional separable part or curves on other kinds of surfaces.

\begin{proposition}\label{prop:normalized-mup-conductor}
Let $X$ be a smooth projective integral surface over an algebraically
closed field $k$ of characteristic $p>0$. Let $Y\subset X$ be a smooth
ample integral curve, and let $M$ be a nontrivial line bundle of order
$p$ with $M|_Y\simeq\OO_Y$. Choose a trivialization of $M^p$ and let
$T\to X$ be its $\mu_p$-torsor. Let $N$ be the normalization of $T$ and
$\pi:N\to X$ the resulting finite map.

The logarithmic differential of $M$ has order $pa$ along $Y$ for an
integer $a\geq1$, with nonzero tangential leading coefficient. The full
inverse image $D=\pi^*Y$ is a reduced integral ample Cartier divisor.
At its generic point the ramification index is one, the residue degree
is $p$, and the conductor of $T$ in $N$ has coefficient $a(p-1)$.
The normalization $\widetilde D$ of $D$ has genus $g(Y)$.

Let $\mathfrak c\subset\OO_N$ be the conductor ideal, put $c=a(p-1)$,
and set $\mathcal J=\mathfrak c\otimes\OO_N(cD)$.
Then $\mathcal J$ is a reflexive ideal in $\OO_N$, and its restriction
embeds in $\OO_D$ with finite-length quotient. Write $\ell$ for that
length and put $\delta=p_a(D)-g(Y)$. Without any smoothness or
Gorenstein assumption on $N$, one has
\[
 2g(Y)-2=ap\,Y^2+\frac{\ell+2\delta}{p-1},
 \qquad \ell\geq0,\quad\delta\geq0.
\]
If $2g(Y)-2=ap\,Y^2$, then $D$ is smooth, $N$ is regular near $D$,
$\mathcal J=\OO_N$, and $N$ is Gorenstein. If $N$ is singular at a
point of $D$, then $2g(Y)-2-apY^2\geq2/(p-1)$.
\end{proposition}
\begin{proof}
On a trivialization of $M$, its chosen $p$th-power trivialization is a
unit $u$, and the torsor has equation $z^p=u$. The forms $du/u$ agree
on overlaps because the transition factors are $p$th powers. They give
a global closed form $\omega$ on $X$. It is nonzero: if a rational
trivialization gave $du=0$, then $u$ would be a $p$th power in $k(X)$,
since $k$ is perfect. The corresponding rational section of $M$ would
have divisor zero, because its $p$th power does. Normality of $X$ would
then trivialize $M$. The same nonvanishing makes $z^p-u$ irreducible
over $k(X)$. Since the torsor algebra is locally free over $X$, it
embeds in its generic field and $T$ is integral.

The image of $\omega|_Y$ in $\Omega^1_{Y/k}$ is zero. To see this,
use a trivialization of $M|_Y$; its $p$th power differs from the chosen
one by a global unit on the proper integral curve $Y$, hence by an
element of $k^*$. The conormal sequence and
$H^0(Y,\OO_Y(-Y))=0$ now give $\omega|_Y=0$.
Let $m\geq1$ be its exact order along $Y$. The nonzero leading section
lies in $\Omega^1_{X/k}|_Y\otimes\OO_Y(-mY)$. Its tangential image
is nonzero, since the kernel on global sections is
$H^0(Y,\OO_Y(-(m+1)Y))=0$. Near the generic point, write the leading
term as $t^m(A\,dt+\eta)$, with $t=0$ defining $Y$ and $\eta$ tangential.
The coefficient of $t^{m-1}$ in $d\omega=0$ is
$m\,dt\wedge\eta$. Since $\eta$ is nonzero modulo $t$, we have
$p\mid m$. Thus $m=pa$.

Put $K=k(Y)$. Complete the local ring at its generic point, choosing a
coefficient field compatible with $k$, to obtain $R=K[[t]]$. In the
completed differential module, write
\[
 u=\sum_{i\geq0}u_it^i.
\]
The equality $\operatorname{ord}_t(du/u)=pa$ implies, for $i<pa$,
that $u_i=0$ if $p\nmid i$ and that $du_i=0$ if $p\mid i$.
The kernel of $d:K\to\Omega^1_{K/k}$ is $K^p$. Choose
$v_j^p=u_{pj}$ for $j<a$, and set $v(t)=\sum_{j<a}v_jt^j$. Then
\[
 u-v(t)^p=t^{pa}b(t),\qquad db(0)\ne0.
\]
In particular $b(0)\notin K^p$. The element
$w=(z-v(t))/t^a$ satisfies $w^p=b(t)$, and
\[
 B=R[w]/(w^p-b(t))
\]
is a one-dimensional local ring with maximal ideal $(t)$ and residue
field $K(b(0)^{1/p})$. It is therefore a discrete valuation ring, and
is the normalization of $A=R[z]/(z^p-u)$ in their common fraction field.
This proves ramification index one and residue degree $p$. Normalization
commutes with this completion by excellence, so these are the invariants
of $N$ along $D$ as well.

As $R$-modules the two rings are
\[
 B=\bigoplus_{i=0}^{p-1}Rw^i,
 \qquad A=\bigoplus_{i=0}^{p-1}Rt^{ai}w^i.
\]
The ideal $t^{a(p-1)}B$ is contained in $A$, hence in its conductor.
Conversely, let $c=\sum c_iw^i$ belong to the conductor. For each $i$,
the product $cw^{p-1-i}$ belongs to $A$, and the coefficient of
$w^{p-1}$ in that product is exactly $c_i$. No other exponent reduces
to $p-1$ using $w^p=b(t)$. Thus each $c_i$ is divisible by
$t^{a(p-1)}$, which proves that the conductor is exactly
$t^{a(p-1)}B$.

The extension of function fields is purely inseparable of degree $p$,
so $\pi$ is radicial. More explicitly, for an integral element $b$ of
the normalization over a normal affine open of $X$, its $p$th power
belongs to that open's fraction field and is integral over its ring;
normality places it in the ring. Thus every prime has a unique prime
above it. The local equation of $Y$ defines a Cartier divisor on $N$;
the calculation shows that this divisor has one component of generic
multiplicity one. Normality gives Serre's condition $S_2$, so its
Cartier quotient is $S_1$. An $S_1$ scheme that is generically reduced
is reduced. Hence the entire inverse image is the integral divisor
$D$. It is ample because it is the finite pullback of an ample divisor.
The smooth projective normalization $\widetilde D$ maps purely
inseparably with degree $p$ to $Y$, so the function-field argument in
the preceding proposition gives $g(\widetilde D)=g(Y)$.

We now compute the correction without resolving $N$. Let $\tau:T\to X$
and $\nu:N\to T$ be the two finite maps. Projection to degree zero
in the cyclic algebra $\mathcal A=\bigoplus_{i\in\mathbf Z/p}M^i$
gives a perfect multiplication pairing: each invertible summand pairs
with its opposite. Thus
$\operatorname{Hom}_{\OO_X}(\mathcal A,\OO_X)\simeq\mathcal A$
as $\mathcal A$-modules. Finite duality gives
$\omega_T\simeq\tau^*\omega_X$. This pairing is not the field trace,
which vanishes for the inseparable function-field extension.

The normal surface $N$ is Cohen--Macaulay. Finite duality for $\nu$
gives
\[
 \omega_N\simeq
 \operatorname{Hom}_{\OO_T}(\nu_*\OO_N,\omega_T)
 \simeq\mathfrak c\otimes\pi^*\omega_X;
\]
see \cite[Tags 0AX2 and 0FVZ]{Stacks}. The second identification
can be checked in a finite birational inclusion of domains $A\subset B$:
every $A$-linear map $B\to A$ is multiplication by an element of their
common fraction field, and its allowed multipliers are exactly the
conductor. The module $\omega_N$ is Cohen--Macaulay of rank one,
hence torsion-free and $S_2$. It is reflexive on the normal surface
\cite[Tag 0AY6]{Stacks}. The same holds for $\mathfrak c$ and
$\mathcal J$.

At the generic point of $D$, the conductor valuation makes
$\mathcal J=\OO_N$; at every other prime divisor its valuation is
nonnegative. Normality therefore gives
$\mathcal J\subset\OO_N$; equivalently, apply the codimension-one
extension criterion \cite[Tag 0AY7]{Stacks}.
At a closed point of $D$, a local equation $t$ is a nonzerodivisor
on the depth-two module $\mathcal J$. The quotient $\mathcal J/t\mathcal J$
has depth one and is torsion-free on the integral curve $D$.
Its map to $\OO_D$ is generically an isomorphism. Any kernel would
have finite support and be torsion, so the kernel is zero. Hence
\[
 0\longrightarrow\mathcal J|_D\longrightarrow\OO_D
 \longrightarrow\mathcal Q\longrightarrow0,
 \qquad\ell=\operatorname{length}_k\mathcal Q\geq0.
\]
This depth argument is what justifies restriction of the possibly
noninvertible conductor ideal.

Cartier-divisor duality \cite[Tag 0AA4]{Stacks}, applied to
$\omega_N[2]$, gives
$\omega_D\simeq(\omega_N\otimes\OO_N(D))|_D$.
There is no higher derived restriction, because $t$ is regular on
$\omega_N$. Define the invertible sheaf
\[
 H=(\pi^*\omega_X\otimes\OO_N((1-c)D))|_D.
\]
Then $\omega_D=H\otimes\mathcal J|_D$, and tensoring the preceding
sequence gives $0\to\omega_D\to H\to H\otimes\mathcal Q\to0$.
Serre duality on the proper Cohen--Macaulay curve gives
$\chi(\omega_D)=-\chi(\OO_D)$
\cite[Tag 0FW0]{Stacks}. The Euler-characteristic definition of degree
therefore yields
\[
 2p_a(D)-2=\deg H-\ell.
\]
Degree multiplies by the degree $p$ of $D\to Y$, even when $D$ is
singular \cite[Tag 0AYZ]{Stacks}. Thus
\[
 D^2=pY^2,\qquad
 \deg H=pK_X\cdot Y+(1-c)pY^2
       =p(2g(Y)-2)-cpY^2.
\]
The normalization sequence gives $p_a(D)=g(Y)+\delta$ with
$\delta\geq0$. Substitution proves the identity.

In the equality case, $\ell=\delta=0$. Vanishing of $\delta$ makes
$D$ normal, hence smooth. At a closed point of $D$, the regular
one-dimensional quotient $\OO_{N,x}/(t)$ has a principal maximal
ideal. Lifting its generator and adjoining $t$ gives two generators
for the maximal ideal of the two-dimensional ring $\OO_{N,x}$;
that ring is regular. The regular locus is open, so $N$ is regular
near $D$.

Vanishing of $\ell$ and Nakayama's lemma make
$\mathcal J=\OO_N$ near $D$. If these reflexive rank-one sheaves
differed anywhere, they would differ along a prime divisor. That
projective curve would be disjoint from $D$, so the canonical section
of $\OO_N(D)$ would trivialize its restriction there. This contradicts
positivity of the degree of an ample line bundle on a curve. Hence
$\mathcal J=\OO_N$ globally. The displayed description of $\omega_N$
now makes it invertible; the Cohen--Macaulay surface is Gorenstein
\cite[Tag 0BFQ]{Stacks}. Finally, a singular point of $N$ on $D$
forces $D$ to be singular by the same local argument, so $\delta\geq1$.
The strict defect bound follows from the identity.
\end{proof}

For the smallest local example, take $p=2$, $K=k(v)$,
$R=K[[t]]$, and $z^2=1+t^2v$. Then $w=(z-1)/t$ satisfies $w^2=v$,
the normalized ring is $B=R[w]$, and $A=R\oplus tRw$ has conductor
$tB$. The uniformizer is still $t$, while the residue field has doubled
its purely inseparable degree. This is a local example of the
normalization calculation, not a global Picard-kernel construction.

If $N$ is smooth near $D$ and its conductor divisor there is
$a(p-1)D+E$, the length $\ell$ equals $E\cdot D$. The module formula
proves the same nonnegative correction without that smoothness
hypothesis. Equality still permits isolated Gorenstein singularities
away from $D$. It does not justify replacing the normal surface by a
smooth one while preserving the ample divisor. On a resolution the
total pullback has zero intersection with exceptional curves, so the
smooth ample-curve obstruction of Proposition~\ref{prop:tangent-ample}
does not apply to that pullback.

\subsection{The group action in the equality case}
Equality in the conductor identity does more than control the canonical
sheaf. It makes the normalized cover a torsor under the same deformation
of $\mu_p$ that appeared in Corollary~\ref{cor:smooth-torsor}.
The group scheme is classical: locally it is the characteristic-$p$
Kummer kernel for the deformed multiplicative group
$\mathcal G^\lambda$ of M\'ezard, Romagny and Tossici
\cite[\S2.2, p.~167, and Lemma~6.1, p.~188]{MRT}.
The assertion here is that the numerical equality forces this action
to extend to the normalization and to be free.

\begin{corollary}\label{cor:normalized-equality-torsor}
In the notation and hypotheses of
Proposition~\ref{prop:normalized-mup-conductor}, suppose
$2g(Y)-2=apY^2$. Let $L=\OO_X(aY)$ and let $s_a$ be its canonical
section. Define $G_{s_a}$ by the formula in
Corollary~\ref{cor:smooth-torsor}, with $L$ and $s_a$ in place of
$\OO_X(Y)$ and $s$. Then $\pi:N\to X$ is a $G_{s_a}$-torsor.
Its infinitesimal action is a homomorphism
\[
 \OO_N(-aD)\longrightarrow T_N
\]
that is nonzero on every fibre, is tangent to $D$, and restricts to
an isomorphism $\OO_D(-aD)\simeq T_D$.
If $a=1$, the surface $N$ has at least one singular point outside $D$.
\end{corollary}
\begin{proof}
The preceding proposition gives the equality of ideal sheaves
$\mathfrak c=\OO_N(-a(p-1)D)$. In particular, $N=T$ away from $D$.
The normal surface $N$ is Cohen--Macaulay and $X$ is regular.
The dimension formula for their finite dominant map and miracle
flatness show that $\pi$ is finite flat of rank $p$
\cite[Tag 00R4]{Stacks}.

The original $\mu_p$-action gives an $X$-derivation $\delta$ of the
function field of $N$. In a local presentation $z^p=u$, it is
characterized by $\delta(z)=z$. It is independent of the chosen
frame of $M$, since changing that frame multiplies $z$ by a function
from $X$. It satisfies $\delta^p=\delta$.

Choose a local equation $t$ for $Y$ and put $\lambda=t^a$ and
$\partial=\lambda\delta$. At the generic point of $D$, the completed
normalization calculation in the proposition writes
$w=(z-v(t))/t^a$. Since $v(t)$ belongs to the completed base ring,
\[
 \partial(t)=0,\qquad \partial(w)=z.
\]
The right side is a unit. Thus $\partial$ preserves the completed
discrete valuation ring. It also preserves the uncompleted one:
its value on an element belongs to the original fraction field and
has nonnegative valuation after completion. The residue derivation is
nonzero as well: a lift to the uncompleted ring of the residue class of
$w$ differs from $w$ by $t$ times the completed ring, and its derivative
therefore has the same nonzero residue as $z$. At every other prime
divisor, $N=T$ and the original action is regular. On an affine open
of the normal integral scheme $N$, its ring is the intersection of
its local rings at height-one primes inside its fraction field.
It follows that $\partial$ is a regular derivation everywhere.

If $\lambda_j=g_{ij}^{-1}\lambda_i$ for two frames of $L$, then
$\partial_j=g_{ij}^{-1}\partial_i$. These derivations therefore
give the displayed homomorphism from $\OO_N(-aD)$. They annihilate
pullbacks from $X$, so their restriction to the smooth curve $D$
is tangent to $D$. The resulting map to $T_D$ is nonzero at its
generic point, by $\partial(w)=z$. Its source and target have the
same degree:
\[
 -aD^2=-apY^2=2-2g(D).
\]
A nonzero map between line bundles of the same degree on a smooth
proper curve is an isomorphism. Hence the infinitesimal action has
no zero on $D$. Away from $D$, its value on the unit $z$ is the unit
$\lambda z$, so it has no zero there either.

We now construct the full group action, rather than only its
derivation. On the chart under consideration, set
\[
 P_0=1,\qquad
 P_i=\frac{\partial(\partial-\lambda)\cdots
 (\partial-(i-1)\lambda)}{i!}\quad(1\leq i<p),
 \qquad
 \rho(b)=\sum_{i=0}^{p-1}P_i(b)u^i.
\]
Here $1$ in $P_0$ denotes the identity operator, and the integers
$i!$ are units in $k$. All these operators are regular, and
$\partial(\lambda)=0$. Where $\lambda$ is invertible, this is the
original $\mu_p$-coaction written through $1+\lambda u$.
Indeed, on an element of weight $j$, for $0\leq j<p$, it is
$b\mapsto b(1+\lambda u)^j$, since
$P_i(b)=\lambda^i\binom ji b$.

Multiplicativity, the identity rule, and coassociativity therefore
hold on this dense open. They hold everywhere because their target
modules, obtained by adjoining one or two variables with $p$th
power zero, are free over the integral local rings of $N$.
On overlaps, each operator $P_i$ scales by $g_{hj}^{-i}$ when the
frame changes by $g_{hj}$; this follows from the fact that the
derivations annihilate $g_{hj}$. Together with $u_j=g_{hj}u_h$, it proves that
the coactions glue. This gives an action of $G_{s_a}$ on $N$ over $X$.

It remains to prove the torsor property. Let $n$ be a closed point
of $N$, lying over $x\in X$. All residue fields are $k$. Since
$\pi$ is radicial and finite flat of rank $p$, its fibre at $x$
is a local scheme of length $p$. The group fibre has coordinate
ring $k[u]/(u^p)$. Nonvanishing of $\partial$ at $n$ gives a local
function whose pullback by the orbit map, after subtracting its
constant term, is $cu+O(u^2)$ with $c\ne0$. Such an element
generates $k[u]/(u^p)$ as an algebra: a truncated compositional
inverse recovers $u$. The orbit map is therefore a closed immersion
between schemes of equal length, hence an isomorphism.

The canonical map $N\times_XG_{s_a}\to N\times_XN$ is a map
between finite locally free $N$-schemes of the same rank $p$.
It is an isomorphism on every closed fibre by the preceding argument.
Its map of finite locally free algebras is consequently an
isomorphism, by its determinant and Nakayama's lemma. Checking
closed points suffices because $N$ is of finite type over $k$.
The finite flat map $\pi$ is surjective, so it is faithfully flat;
the action thus makes it a torsor.

Finally, if $a=1$ and $N$ were smooth, the everywhere nonzero
infinitesimal action would be a line subbundle of $T_N$ with
locally free quotient, tangent to the smooth ample curve $D$.
This contradicts Proposition~\ref{prop:tangent-ample}.
The preceding proposition already gives regularity near $D$,
so a singular point must lie outside $D$.
\end{proof}

The possibility of a singular torsor over a smooth surface is already
visible locally. On the open surface where $1+xy$ is invertible,
the equation $z^p=1+xy$ defines a $\mu_p$-torsor, because $z$ is a
unit. At $z=1$, putting $w=z-1$ gives $xy=w^p$. This is an integral
two-dimensional hypersurface, regular away from its origin and hence
normal by Serre's criterion. It is Gorenstein but singular at the
origin. Thus the torsor property does not eliminate the isolated
singularities allowed by the equality case.

\subsection{The isolated singularities and the tangent bundle}
The remaining singularities can be measured on the smooth base. Dividing
the logarithmic form by its exact vanishing along $Y$ gives a regular
section of a rank-two vector bundle. Its zero scheme detects precisely
the singular locus of the normalized cover and has a computable length.

\begin{proposition}\label{prop:normalized-mup-singularity-defect}
Keep the notation of Proposition~\ref{prop:normalized-mup-conductor},
and suppose that
\[
 2g(Y)-2=apY^2.
\]
Let $\omega$ be the logarithmic differential of $M$, put $m=pa$,
and let $s_Y$ be the canonical section of $\OO_X(Y)$.
The section
\[
 \theta=\omega/s_Y^m\in
 H^0\bigl(X,\Omega^1_{X/k}(-mY)\bigr)
\]
has a finite zero scheme $Z$ disjoint from $Y$, and
\begin{equation}\label{eq:normalized-mup-zero-length}
 \operatorname{length}_k Z=c_2(X)+apY^2=c_2(X)+aD^2.
\end{equation}
The singular subscheme of $N$ defined by
$\operatorname{Fitt}_2(\Omega^1_{N/k})$ is $N\times_X Z$.
Its length is $p\operatorname{length}_k Z$, and $N$ is smooth if and
only if $Z$ is empty.

The tangent sheaf $T_N$ is a vector bundle of rank two, even when $N$
is singular. The infinitesimal action of
Corollary~\ref{cor:normalized-equality-torsor} gives an exact sequence of vector bundles
\begin{equation}\label{eq:normalized-mup-tangent-sequence}
 0\longrightarrow\OO_N(-aD)\longrightarrow T_N
 \longrightarrow\omega_N^{-1}(aD)\longrightarrow0,
\end{equation}
and
\begin{equation}\label{eq:normalized-mup-tangent-chern}
 \deg c_2(T_N)=-aD^2.
\end{equation}
If $a=1$, then $Z$ is nonempty, so $c_2(X)+pY^2\geq1$.
\end{proposition}

\begin{proof}
The equality conclusions of
Proposition~\ref{prop:normalized-mup-conductor} give that $D$ is
smooth, $N$ is regular near $D$, and
\[
 \mathfrak c=\OO_N\bigl(-a(p-1)D\bigr),\qquad
 \omega_N\simeq\pi^*\omega_X\otimes
                \OO_N\bigl(-a(p-1)D\bigr).
\]
In particular the normalization is an isomorphism away from $D$:
$N\setminus D=T\setminus\tau^{-1}Y$. The map $\pi$ is finite flat
of rank $p$. Indeed, its source is a Cohen--Macaulay normal surface,
its target is regular, and its local fibres have dimension zero;
the local dimensions agree because the map is finite and radicial.
Miracle flatness applies \cite[Tag 00R4]{Stacks}.

The exact order $m$ of $\omega$ makes $\theta$ regular. Its
tangential component along $Y$ is a nonzero section of
\[
 \Omega^1_{Y/k}\otimes\OO_Y(-mY).
\]
This line bundle has degree zero by the equality hypothesis, so the
section is nowhere zero. Thus $\theta$ has no zero on $Y$.

Away from $Y$, a trivialization of $M$ gives
$T=\Spec R[z]/(z^p-u)$ with $u\in R^*$, and $\theta$ is a unit
multiple of $du$. Choose an \'etale local basis $dx,dy$ of
$\Omega^1_{R/k}$ and write $du=\alpha\,dx+\beta\,dy$.
The differential module of this hypersurface has three generators
$dx,dy,dz$ and the single relation
$\alpha\,dx+\beta\,dy=0$. Its second Fitting ideal is therefore
$(\alpha,\beta)$. The hypersurface is smooth precisely where this
ideal is the unit ideal.

On this open set $T=N$ is normal, so its singular locus has no
codimension-one component. Finiteness of $\pi$ implies that the
zero scheme of $\theta$ has no curve away from $Y$. Since it also
avoids $Y$, the scheme $Z$ is zero-dimensional, hence finite. At a
point of $Z$, its two local equations form a regular sequence in
the regular local ring of the smooth surface $X$. The top Chern
class of a regular section now gives
\[
 \operatorname{length}_k Z
 =\deg c_2\bigl(\Omega^1_{X/k}(-mY)\bigr);
\]
all the dimension and regular-sequence hypotheses of
\cite[Tag 0FA9]{Stacks} hold here. The rank-two twist formula
\cite[Tag 02UD]{Stacks} gives
\[
 \deg c_2\bigl(\Omega^1_{X/k}(-mY)\bigr)
 =c_2(X)-mK_X\cdot Y+m^2Y^2.
\]
Adjunction and equality imply $K_X\cdot Y=(m-1)Y^2$, proving
\eqref{eq:normalized-mup-zero-length}.

The local Fitting-ideal computation also proves
$\operatorname{Fitt}_2(\Omega^1_{N/k})=\mathcal I_Z\OO_N$
off $D$. Near $D$ both ideals are the unit ideal, because $N$ is
smooth there and $Z$ avoids $Y$. The equality therefore holds
globally. Flatness of rank $p$ gives the asserted length. The same
local hypersurface calculation, together with smoothness near $D$,
shows that $N$ is smooth exactly when $Z$ is empty.

We next check local freeness of the tangent sheaf at a singular
point. In the above presentation put $B=R[z]/(z^p-u)$.
The elements $\alpha,\beta$ form a regular sequence in $R$,
because their zero scheme is finite, and remain a regular sequence
in $B$ by flatness. Thus
\[
 \operatorname{Der}_k(B,B)
 =\ker\bigl(B^3\longrightarrow B,\ (r,s,t)\longmapsto
             \alpha r+\beta s\bigr)
\]
is free with basis $(-\beta,\alpha,0)$ and $(0,0,1)$: the first
generator is the Koszul syzygy of the regular sequence. At every
other point $N$ is smooth. Hence $T_N$ is a rank-two vector bundle
on all of $N$.

Corollary~\ref{cor:normalized-equality-torsor} gives the map
$\OO_N(-aD)\to T_N$, nonzero in every tangent fibre and tangent
to $D$. Since $T_N$ is now known to be locally free, this is a
line subbundle with invertible quotient.

On the regular locus, $\det T_N=\omega_N^{-1}$. Both sheaves are
invertible on all of $N$, and normality extends this identification
across the finite complement. The quotient is therefore
$\omega_N^{-1}(aD)$, proving
\eqref{eq:normalized-mup-tangent-sequence}. Adjunction along $D$
gives $K_N\cdot D=(a-1)D^2$. The Whitney formula then yields
\[
 \deg c_2(T_N)=(-aD)\cdot(-K_N+aD)=-aD^2,
\]
as claimed.

Finally, if $a=1$ and $Z$ were empty, $N$ would be a smooth
projective surface with the line subbundle $\OO_N(-D)\subset T_N$
tangent to its smooth ample curve $D$. This contradicts
Proposition~\ref{prop:tangent-ample}. Thus $Z$ has positive integral
length in this case.
\end{proof}

The length in this proposition is a Fitting-scheme length. It is not
a claimed Milnor number or a discrepancy correction on a resolution.
In particular, the negative second Chern number of the vector bundle
$T_N$ cannot be substituted for the second Chern number of a smooth
resolution. The local model $xy=w^p$ already has a locally free
tangent sheaf and a singular subscheme of length $p$.

\begin{corollary}\label{cor:normalized-mup-plurigenera}
In Proposition~\ref{prop:normalized-mup-singularity-defect}, suppose
$a=1$. Let $f:S\to N$ be a smooth projective resolution that is an
isomorphism over the regular locus of $N$. For every integer $r>0$
there is an injection
\[
 H^0(S,\omega_S^{\otimes r})\longrightarrow
 H^0(N,\omega_N^{\otimes r}),
 \qquad h^0(N,\omega_N^{\otimes r})\leq1.
\]
Every nonzero section of $\omega_N^{\otimes r}$ is nowhere
vanishing. Thus $\kappa(S)\leq0$, and $K_N$ is torsion if some
positive plurigenus of $S$ is nonzero.
\end{corollary}

\begin{proof}
Here $K_N\cdot D=0$, and $D$ is ample. A nonzero section of the
invertible sheaf $\omega_N^{\otimes r}$ has an effective Cartier
zero divisor $E$. If $E$ were nonempty, ampleness would give
$E\cdot D>0$, whereas $E\cdot D=rK_N\cdot D=0$.
Nor can the section have only an isolated zero: a nonzero nonunit
in a noetherian normal local domain has a height-one zero.
The section therefore trivializes $\omega_N^{\otimes r}$.
Since $N$ is proper and integral over the algebraically closed
field, $H^0(N,\OO_N)=k$. It follows that the indicated space of
sections has dimension at most one.

Restrict a pluricanonical section on $S$ to the common regular
open set. The complement of that open set in $N$ is finite.
Normality extends the section uniquely to the invertible sheaf
$\omega_N^{\otimes r}$ on $N$ \cite[Tag 0AY7]{Stacks}.
This gives the injection, which is injective because $S$ is
integral. All positive plurigenera of $S$ are consequently at
most one. If one is nonzero, the corresponding power of
$\omega_N$ is trivial by the first paragraph. The claimed
Kodaira-dimension bound follows.
\end{proof}

For the smallest equality case, Ekedahl's foliation theorem
restricts the surface further. Following his terminology, a surface
is \emph{almost ruled} if it admits a dominant purely inseparable
rational map of degree $p$ from a ruled surface
\cite[p.~108]{EkedahlCanonical}.

\begin{proposition}[After Ekedahl]\label{prop:normalized-mup-ekedahl}
In Proposition~\ref{prop:normalized-mup-conductor}, suppose $a=1$
and $2g(Y)-2=pY^2$. Then either $X$ is almost ruled, or $p=2$ and
\[
 \omega_X\text{ is ample},\qquad
 \chi(\OO_X)=1,\qquad
 \omega_X^{\otimes2}\simeq\OO_X(2Y).
\]
In the latter case $N$ has only rational double points, its minimal
resolution is a K3 surface, and
\[
 1\leq Y^2=K_X^2\leq9,\qquad g(Y)=Y^2+1.
\]
In this latter case the zero scheme $Z$ of
Proposition~\ref{prop:normalized-mup-singularity-defect} has length
$12+Y^2$, and the Fitting singular subscheme of $N$ has length
$2(12+Y^2)$. In particular, when $p>2$ the stated equality forces
$X$ to be almost ruled.
\end{proposition}

\begin{proof}
The wedge pairing for the rank-two cotangent bundle identifies
$\Omega^1_{X/k}$ with $T_X\otimes\omega_X$. Thus the nonzero
section $\theta\in H^0(X,\Omega^1_{X/k}(-pY))$ gives an injection
\[
 \mathcal L:=\OO_X(pY)\otimes\omega_X^{-1}\longrightarrow T_X.
\]
Its image is saturated. Indeed, in smooth local coordinates write
$\theta=A\,dx+B\,dy$. Its zero scheme is finite, so at a zero the
pair $A,B$ is a regular sequence. The kernel of
$(r,s)\mapsto Ar+Bs$ is generated by $(-B,A)$, and its image is
the torsion-free ideal $(A,B)$. Globally this gives
\[
 0\longrightarrow\mathcal L\longrightarrow T_X
 \longrightarrow\mathcal I_Z\otimes\OO_X(-pY)
 \longrightarrow0.
\]
On a Kummer trivialization the form is a nonzero rational multiple
of $du$, where the cover has equation $z^p=u$. Its annihilator
consists of derivations killing $u$, hence is closed under brackets
and $p$th powers. Saturation extends these properties to the whole
sheaf. This is a possibly singular rank-one foliation in the sense
of \cite[pp.~104--105]{EkedahlCanonical}.

Adjunction and the equality hypothesis give
\[
 K_X\cdot Y=(p-1)Y^2,\qquad
 \mathcal L\cdot Y=Y^2>0,
 \qquad (K_X-(p-1)\mathcal L)\cdot Y=0.
\]
Assume that $X$ is not almost ruled. It is then not birationally
ruled: a ruled function field $k(C)(t)$ has the purely inseparable
degree-$p$ extension $k(C)^{1/p}(t)$, which is again ruled.
We may therefore apply
\cite[I, Theorem~2.4(ii), pp.~109--115]{EkedahlCanonical}, with
the injection just constructed and the ample test line bundle
$\OO_X(Y)$.

In alternative~(a) of that theorem, $X$ is minimal of Kodaira
dimension zero or one and
\[
 \omega_X^{\otimes12p}\simeq
 \mathcal L^{\otimes12p(p-1)}.
\]
Substitution gives
$\omega_X^{\otimes12p^2}\simeq
\OO_X(12p^2(p-1)Y)$, which makes $\omega_X$ ample and is
impossible in those Kodaira dimensions. Alternative~(b) makes
$\mathcal L$ trivial, contrary to $\mathcal L\cdot Y>0$.
Alternative~(c) has $p=2$, Kodaira dimension one, and
$\omega_X^{\otimes2}\simeq\mathcal L^{\otimes2}$.
It would give $\omega_X^{\otimes4}\simeq\OO_X(4Y)$,
again a contradiction. Alternative~(d) consequently holds:
$p=2$, $\chi(\OO_X)=1$ and $\mathcal L\simeq\omega_X$.
This proves the asserted square identity and ampleness. It also
shows that $X$ is minimal, since a $(-1)$-curve would have negative
intersection with $K_X$.

We identify the foliation cover in that theorem with $N$.
Put $K=k(X)$ and $E=K(z)$, with $z^p=u$. The nonzero form gives
$u\notin K^p$. Complete $u$ to a $p$-basis $u,v$ of $K$ over
$K^p$. Expansion in the basis $u^iv^j$, for $0\leq i,j<p$,
shows that the invariant field of the derivations killing $u$ is
$F=K^p(u)$. On the other hand,
\[
 E^p=K^p(u)=F,\qquad E=F^{1/p}.
\]
Ekedahl's convention has relative Frobenius $X^{(1)}\to X$;
his foliation cover $(X/\mathcal L)^{(1)}\to X$ has field
$F^{1/p}$ \cite[p.~104]{EkedahlCanonical}.
The identification is over $X$, not merely birational. For a
normal affine ring $R$ on $X$, the invariant ring is $C=R\cap F$.
It is normal and contains $R^p$. The ring $C^{1/p}$ contains $R$,
is a submodule of the finite $R$-module $R^{1/p}$, and is normal
with fraction field $E$. Hence it is the integral closure of $R$
in $E$. These descriptions glue and identify that cover with $N$.

For completeness we spell out the part of the source's proof
that gives the assertion about this cover. Let $\tau:S\to N$
be the minimal resolution and put $f=\pi\tau$. The canonical
formula in \cite[I, Proposition~2.1, p.~108]{EkedahlCanonical}
and the zero intersection above give $K_S\cdot f^*Y=0$.
Thus $\kappa(S)\leq0$, also by
Corollary~\ref{cor:normalized-mup-plurigenera}. If
$\kappa(S)=-\infty$, the surface $S$ is birationally ruled and
its degree-$p$ purely inseparable map to $X$ makes $X$ almost
ruled. Consequently $\kappa(S)=0$.

The surface $S$ is minimal. Otherwise its canonical class is the
pullback of the numerically trivial canonical class of its minimal
model plus an effective exceptional divisor. The zero intersection
with $f^*Y$ forces every component of that divisor to be contracted
by $f$, hence by $\tau$, because $N\to X$ is finite. This
contradicts minimality of the resolution. The surface $N$ is
singular by Proposition~\ref{prop:normalized-mup-singularity-defect}.
Thus \cite[I, Lemma~2.5, p.~111]{EkedahlCanonical} applies:
$N$ has rational double points, and $S$ is K3 or Enriques.

The Enriques possibility is excluded by Euler characteristic.
A finite flat purely inseparable double cover has zero trace and is an
$\alpha_{\mathcal A}$-torsor for an invertible sheaf $\mathcal A$,
and its algebra fits into
\[
 0\longrightarrow\OO_X\longrightarrow\pi_*\OO_N
 \longrightarrow\mathcal A^{-1}\longrightarrow0
 \quad\text{\cite[I, Proposition~1.11, p.~107]{EkedahlCanonical}}.
\]
The associated foliation has
$\mathcal L\simeq\mathcal A^{-2}\otimes\omega_X^{-1}$.
Since $\mathcal L\simeq\omega_X$, the bundle
$\mathcal A^{-1}$ is numerically equivalent to $\omega_X$.
Riemann--Roch gives
$\chi(\mathcal A^{-1})=\chi(\omega_X)=\chi(\OO_X)=1$.
Rational double points are rational singularities, so
$\chi(\OO_S)=\chi(\OO_N)=2$. An Enriques surface has Euler
characteristic one. Therefore $S$ is K3.

The purely inseparable morphism $S\to X$ preserves the first
$\ell$-adic Betti number
\cite[Lemma~0.1.5, p.~103]{EkedahlCanonical}, so $b_1(X)=0$.
Poincar\'e duality gives $c_2(X)=2+b_2(X)\geq3$: the ample
divisor $Y$ has nonzero cycle class because its square is positive.
Noether's formula gives $c_2(X)=12-K_X^2$, while the square
identity gives $K_X^2=Y^2>0$. This proves $1\leq Y^2\leq9$.
The genus formula is the assumed equality for $p=2$.
Finally, \eqref{eq:normalized-mup-zero-length} gives
$\operatorname{length}_k Z=12+Y^2$, and the rank-two flat cover
doubles that length on $N$.
\end{proof}

This proposition applies Ekedahl's published theorem. The
almost-ruled alternative and the characteristic-two K3 case
remain open for the Picard restriction problem. In particular,
a purely inseparable map from a ruled surface does not assert
that $X$ itself is birationally ruled. The strict inequality
$2g(Y)-2>apY^2$ and the cases $a>1$ also remain outside this
reduction.
